\section*{第 \ref{ch:b1:magnetostatics} 章——磁静力学：场、力与偶极子}

\begin{solution}{exo:b1:magnetostatics:1}
$B = \mu_0I/2\pi r = 2 \times 10^{-7} \times 10/0.01 = \qty{2e-4}{T}$。若取 \qty{1}{T}：
$I = 2\pi rB/\mu_0 = \qty{5e4}{A}$。地磁场所对应的 $r = 2 \times 10^{-7} \times
100/5 \times 10^{-5} = \qty{0.4}{m}$。
\end{solution}

\begin{solution}{exo:b1:magnetostatics:2}
圆心处：$\mu_0NI/2R = 4\pi \times 10^{-7} \times 200/0.1 = \qty{2.5}{mT}$。当 $z = R$ 时：
$\times(R^2/2R^2)^{3/2} = 2^{-3/2}$：\qty{0.89}{mT}。当 $z = 10R$ 时，$\times(1/101)^{3/2}
= 10^{-3}$：\qty{2.5}{\micro T}，已符合偶极子的 $1/z^3$ 规律。
\end{solution}

\begin{solution}{exo:b1:magnetostatics:3}
$B = \mu_0nI = 4\pi \times 10^{-7} \times 1000 \times 5 = \qty{6.3}{mT}$；在端点处，
减半为 \qty{3.1}{mT}。线径 \qty{0.5}{mm} 的导线：每米 $2000$ 匝，故
在 \qty{5}{A} 时为 \qty{12.6}{mT}。
\end{solution}

\begin{solution}{exo:b1:magnetostatics:4}
$F = ILB = 10 \times 1 \times 0.1 = \qty{1.0}{N}$；当 \qty{30}{\degree} 时，$\sin30^\circ$：
\qty{0.5}{N}。电动机：$100 \times \qty{1}{N} \times \qty{0.05}{m} = \qty{5}{N.m}$——
一台小型工业电动机。
\end{solution}

\begin{solution}{exo:b1:magnetostatics:5}
令 $z = d\tan\alpha$：$B = \dfrac{\mu_0I}{4\pi}\displaystyle\int\dfrac{d\,\dd z}{(d^2 + z^2)^{3/2}}
= \dfrac{\mu_0I}{4\pi d}\displaystyle\int_{\alpha_1}^{\alpha_2}\cos\alpha\,\dd\alpha = \dfrac{\mu_0I}{4\pi d}
(\sin\alpha_2 - \sin\alpha_1)$；无限长导线，$\alpha_{1,2} = \mp\pi/2$：$\mu_0I/2\pi d$。
正方形：每边位于 $d = a/2$ 处，所张角为 $\pm45^\circ$，给出 $\mu_0I\sqrt2/2\pi a$；
四边合计：$B = 2\sqrt2\,\mu_0I/\pi a = 0.90\,\mu_0I/a$（作为对比，一个“半径”
为 $a/2$ 的圆将给出 $\mu_0I/a$）。
\end{solution}

\begin{solution}{exo:b1:magnetostatics:6}
$B(z) = f(z - R/2) + f(z + R/2)$，其中 $f(u) = \mu_0NIR^2/2(R^2 + u^2)^{3/2}$。
中点处：$2f(R/2) = \mu_0NIR^2/(5R^2/4)^{3/2} = (4/5)^{3/2}\mu_0NI/R$。$f'$ 为
奇函数，故 $B'(0) = f'(-R/2) + f'(R/2) = 0$；$f''(u) \propto (4u^2 - R^2)(R^2 +
u^2)^{-7/2}$ 在 $u = R/2$ 处为零，故 $B''(0) = 2f''(R/2) = 0$，场直至
三阶均匀。代入数值：$0.716 \times 4\pi \times 10^{-7} \times 100/0.15 =
\qty{0.60}{mT}$。
\end{solution}

\begin{solution}{exo:b1:magnetostatics:7}
表面处：$\mu_0I/2\pi a = 2 \times 10^{-7} \times 100/0.002 = \qty{10}{mT}$；在
\qty{1}{mm} 处（内部）$\times r/a$：\qty{5}{mT}；在 \qty{1}{cm} 处：\qty{2}{mT}。$B$ 在
表面以内线性增大，随后按 $1/r$ 减小。同轴情形下，护套外
所包围的电流为零，$B = 0$。
\end{solution}

\begin{solution}{exo:b1:magnetostatics:8}
$2 \times 10^{-7}$~N/m。母线：$2 \times 10^{-7} \times 10^8/0.1 = \qty{200}{N/m}$；
在 \qty{100}{kA} 时：$\qty{2e4}{N/m}$——每米两吨。开关设备的绝缘
支撑件按故障电流而非
正常工作电流来设计。
\end{solution}

\begin{solution}{exo:b1:magnetostatics:9}
$\Gamma = mB_h = \qty{2e-6}{N.m}$；$W = 2mB_h = \qty{4e-6}{J}$；$J\ddot\theta =
-mB_h\theta$：$T = 2\pi\sqrt{J/mB_h} = 2\pi\sqrt{10^{-7}/2 \times 10^{-6}} =
\qty{1.4}{s}$。测定振荡周期，由 $B_h = 4\pi^2J/mT^2$ 求得，前提是 $m$ 与 $J$
均已知（高斯单独测得 $m$，方法是
利用它使另一根指针产生的偏转）。
\end{solution}

\begin{solution}{exo:b1:magnetostatics:10}
(a) $\mu_0IR^2/2z^3 = \mu_0(I\pi R^2)/2\pi z^3 = \mu_0m/2\pi z^3 = (\mu_0/4\pi)\,2m/z^3$：
对应电场中的 $2p/4\pi\varepsilon_0z^3$，只需 $1/\varepsilon_0 \to \mu_0$。(b) $m = 2\pi R_E^3B/
\mu_0 = 2\pi \times 2.58 \times 10^{20} \times 6 \times 10^{-5}/1.26 \times 10^{-6} = \qty{7.7e22}{A.m^2}$。
(c) $I = m/\pi R^2 = 7.7 \times 10^{22}/(\pi \times 9 \times 10^{12}) = \qty{2.7e9}{A}$。(d)
赤道处（$\theta = 90^\circ$）：$B_\theta = \mu_0m/4\pi R_E^3 = \qty{3e-5}{T}$，为
极区值的一半。在 $\theta = 45^\circ$ 处：$\tan(\text{磁倾角}) = B_r/B_\theta = 2\cot\theta = 2$，
磁倾角 $= \qty{63}{\degree}$——在中纬度，磁场
陡直地插入地下。
\end{solution}

\begin{solution}{exo:b1:magnetostatics:11}
$\mu_0nI = 4\pi \times 10^{-7} \times 2000 \times 3 = \qty{7.5}{mT}$。圆心处：$\tan\alpha_1 =
R/(L/2) = 0.2$，$\cos\alpha_1 = -\cos\alpha_2 = 0.981$：$B = 0.981\,\mu_0nI =
\qty{7.4}{mT}$（$98\%$）。端点处：$\cos\alpha_1 = 0$，$\tan\alpha_2 = R/L$，$\cos\alpha_2 =
-0.995$：$B = 0.497\,\mu_0nI = \qty{3.7}{mT}$（$50\%$）。在超出端点的 $x$
处，两端位于同侧：$B = \tfrac12\mu_0nI(\cos\alpha_1 - \cos\alpha_2)
\approx \tfrac12\mu_0nI\,\tfrac{R^2}{2}[1/x^2 - 1/(x + L)^2]$；约为圆心值的 $1\%$，
在 $x \approx \qty{10}{cm}$ 处，相距约半个长度。$B(z)$ 在
长度的大部分近似平顶，在两端降为一半，并
在离端一个长度处降至百分之几。
\end{solution}

\begin{solution}{exo:b1:magnetostatics:12}
(a) 绕组自身的场在穿过绕组时由 $0$ 跃变为 $B$；
作用于其上的场为平均值 $B/2$（薄层不能对
自身施力）。长 $\dd l$、宽 $\dd z$ 的绕组微元所受之力，
携带电流 $K\,\dd z$：$K\,\dd z\,\dd l\,B/2$，沿径向向外（可用 $I\,\dd\vect l
\wedge\vect B$ 检验）：压强 $p = KB/2 = B^2/2\mu_0$，因 $K = nI = B/\mu_0$。
(b) \qty{1}{T}：\qty{4e5}{Pa} = \qty{4}{bar}；\qty{10}{T}：\qty{400}{bar}；\qty{45}{T}：
\qty{8e8}{Pa} = \qty{8000}{bar}，接近钢材的 \qty{1}{GPa}。(c) 绕组
须承受随 $B^2$ 增长的压强：在约 \qty{40}{T} 以上
没有任何材料或冷却方案能够应对——稳态磁体的极限
在于力学（对于超导体，则是其临界场）。
\end{solution}

\begin{solution}{pb:b1:magnetostatics:1}
\textbf{1.} 任一包含轴线的平面都是电流分布的
对称面：$\vect B$ 垂直于该平面，沿周向；
由沿轴与绕轴的不变性，$\vect B = B(r)\,\vect e_\theta$。

\textbf{2.} 半径为 $r$ 的圆周：$2\pi rB = \mu_0I_{\text{包围}}$。芯部：
$I_{\text{包围}} = Ir^2/a^2$，$B = \mu_0Ir/2\pi a^2$；中间区域：$\mu_0I/2\pi r$；
外部（$r > c$）：$I - I = 0$，$B = 0$。

\textbf{3.} $B(a) = 2 \times 10^{-7}/5 \times 10^{-4} = \qty{4.0e-4}{T}$; $B(b) = 2 \times
10^{-7}/2.5 \times 10^{-3} = \qty{8.0e-5}{T}$.

\textbf{4.} 管内没有电流被包围，因此此处 $B = 0$，
而不是线性上升；其他区域不变——圆柱形电流外部的
场只取决于总电流。

\textbf{5.} $I_{\text{包围}} = I - I(r^2 - b^2)/(c^2 - b^2) = I(c^2 - r^2)/(c^2 -
b^2)$：$B = \dfrac{\mu_0I}{2\pi r}\,\dfrac{c^2 - r^2}{c^2 - b^2}$，其中 $\mu_0I/2\pi b$ 为
$b$ 处的值，$0$ 为 $c$ 处的值：连续。草图：线性上升至
\qty{4e-4}{T}（\qty{0.5}{mm} 处），再按 $1/r$ 下降至 \qty{8e-5}{T}（\qty{2.5}{mm} 处），
在 \qty{3}{mm} 处降为零，此后恒为零。

\textbf{6.} 护套中的电流与芯部电流反向平行：相互排斥，
产生向外的压强。表面电流 $K = I/2\pi b = \qty{64}{A/m}$；平均场
$\tfrac12(8 \times 10^{-5} + 0) = \qty{4e-5}{T}$：$p = KB = \qty{2.5e-3}{Pa}$。在
\qty{10}{kA} 时：$\times10^8$，即 \qty{2.5e5}{Pa} = \qty{2.5}{bar}——
脉冲功率电缆正是为承受这一压强而设计的。

\textbf{7.} 外部无磁场：该电缆既不向外辐射，
也不会拾取邻近导线的磁干扰。一对导线：两者的场
近乎相互抵消，剩下 $\approx \mu_0Id/2\pi D^2 = 2 \times 10^{-7} \times 2 \times
10^{-3}/10^{-2} = \qty{4e-8}{T}$，位于 \qty{10}{cm} 处——很小，但同轴
电缆的场为零。

\textbf{8.} $D = \sqrt{h^2 + d^2/4} = \qty{10.0}{m}$：每根导线给出 $\mu_0I/2\pi D =
\qty{1.0e-5}{T}$，方向垂直于导线到该点的连线。两电流
方向相反，平行于两导线连线的分量相互抵消，
竖直分量相加：$B = 2 \times 10^{-5} \times
(d/2)/D = \qty{1.0e-6}{T}$，竖直方向。比地磁场小五十倍，
比指导值低一百倍。

\textbf{9.} 当 $D \gg d$ 时，$B \approx \dfrac{\mu_0I}{2\pi}\Bigl(\dfrac1{D_1} - \dfrac1{D_2}\Bigr)
\sim \dfrac{\mu_0I}{2\pi}\dfrac{d}{D^2}$：两个 $1/D$ 的场在一阶近似下相互抵消，
只剩相对大小为 $d/D$ 的差值——
即一个“双线偶极子”。

\textbf{10.} 自身导线位于 \qty{1}{m} 处：\qty{1.0e-4}{T}，水平方向；另一根
位于 \qty{1.41}{m} 处：\qty{7.1e-5}{T}，方向与水平成 \qty{45}{\degree}，与之方向相反：合场
$(5 \times 10^{-5}, 5 \times 10^{-5})$，即 \qty{7e-5}{T} ——\qty{70}{\micro T}，低于
公众指导值，远低于职业暴露限值。

\textbf{11.} $F/\ell = \mu_0I^2/2\pi d = 2 \times 10^{-7} \times 2.5 \times 10^5 =
\qty{0.05}{N/m}$，斥力；在 \qty{20}{kA} 时为 \qty{80}{N/m}。

\textbf{12.} 正常时为 \qty{15}{N}；短路时为 \qty{24}{kN}，是跨距
\qty{2.9}{kN} 重力的八倍：导线被猛然甩开、摆动，
落回时可能相互碰撞——因此需要间隔器，以及能在
数个周期内断开的断路器。

\textbf{13.} $I\cos\omega t + I\cos(\omega t - 2\pi/3) + I\cos(\omega t + 2\pi/3) =
I\cos\omega t\,(1 + 2\cos2\pi/3) = 0$：三个相量之和为零。
$1/D$ 的场如同同轴电缆一样相互抵消；远场以 $1/D^2$
或更快的速率衰减。

\textbf{14.} 该处导线的场是竖直的，而指南针只对
水平分量有响应：不发生偏转（倾角变化
$10^{-6}/(2 \times 10^{-5})$，即三度，指南针无法显示）。

\textbf{15.} 边 $a$（与轴平行）：径向的 $\vect B$ 与之
垂直，每边受力 $NIaB$，方向垂直于
线圈平面，两边相反：构成力偶。边 $b$：$\dd\vect l$
沿径向，与 $\vect B$ 平行：不受力。

\textbf{16.} $\Gamma = 2 \times NIaB \times b/2 = NIabB = NISB = 0.024\,I$（N\,m，$I$
单位为 A）。径向场始终位于线圈平面内，并与边 $a$ 垂直：
各力始终与平面垂直，其力臂
$b/2$ 永不改变。

\textbf{17.} $C\theta = NISB$：$\theta = 2400\,I$ rad，即 \qty{2.4}{rad/mA}；满量程
\qty{1.2}{rad} 对应 $I = \qty{0.50}{mA}$。

\textbf{18.} 长度 $200 \times 2(0.02 + 0.03) = \qty{20}{m}$，截面 $\pi(5 \times
10^{-5})^2 = \qty{7.9e-9}{m^2}$：$R = 1.7 \times 10^{-8} \times 20/7.9 \times 10^{-9} =
\qty{43}{\ohm}$；$U = 43 \times 5 \times 10^{-4} = \qty{22}{mV}$；$P = UI = \qty{11}{\micro W}$。

\textbf{19.} $J\ddot\theta = -C\theta + NISB\,I$：$\omega_0 = \sqrt{C/J} = \sqrt{10^{-5}/6.7
\times 10^{-8}} = \qty{12}{rad/s}$，$T = \qty{0.51}{s}$。

\textbf{20.} $J\ddot\theta + \alpha\dot\theta + C\theta = NISB\,I$：临界阻尼条件为
$\alpha = 2\sqrt{JC} = 2\sqrt{6.7 \times 10^{-13}} = \qty{1.6e-6}{N.m.s}$；指针此时
以最快速度到达读数位置，且不会超调。

\textbf{21.} 电压表：$R_s = 10/5 \times 10^{-4} - 43 \approx \qty{20}{k\ohm}$。
电流表：分流器在 \qty{22}{mV} 电压下通过 \qty{0.9995}{A} 电流：$R_{\text{分}} =
\qty{22}{m\ohm}$。

\textbf{22.} $m = NIS = 200 \times 5 \times 10^{-4} \times 6 \times 10^{-4} = \qty{6e-5}{A.m^2}$；
$-mB = \qty{-1.2e-5}{J}$；弹簧 $\tfrac12C\theta^2 = \tfrac12 \times 10^{-5} \times 1.44 =
\qty{7.2e-6}{J}$：二者量级相同——
弹簧储存的，正是线圈转动时磁力矩所做的功。

\textbf{23.} 在均匀场中，力矩为 $mB\cos\theta$（平面包含 $\vect B$ 时最大，
法线与之平行时为零）：$C\theta =
NISB\,I\cos\theta$——仅在 $\theta$ 较小时线性，
刻度顶端处被压缩。

\textbf{24.} \qty{50}{Hz}，对比 $\omega_0/2\pi = \qty{2}{Hz}$：线圈无法跟随；它
停在电流的平均值（零）处，并伴有幅度为
$\approx \theta_{\text{静态}}(\omega_0/\omega)^2 = 1.2 \times (12/314)^2 \approx
\qty{2e-3}{rad}$ 的微小抖动，肉眼不可见。动圈式仪表读出的是平均值——
对于交流电，需要配备整流器。

\textbf{25.} 在 \qty{1}{A} 同轴电缆的芯部为 \qty{4e-4}{T}，其外部无磁场
（安培环路定理）；在 \qty{500}{A} 导线下方为 \qty{1}{\micro T}，两导线之间为 \qty{0.05}{N/m}
（毕奥–萨伐尔定律加拉普拉斯力）；仪表为 \qty{2.4}{rad/mA}
（拉普拉斯力偶与弹簧平衡）。
\end{solution}
